\documentclass[11pt,a4paper]{article}
\usepackage[margin=1in]{geometry}
\usepackage{listings}
\usepackage{xcolor}
\usepackage{titlesec}
\usepackage{tcolorbox}
\tcbuselibrary{skins,breakable}

\definecolor{outputbg}{rgb}{0.92,0.96,0.92}
\definecolor{outputborder}{rgb}{0.3,0.6,0.3}

\lstset{
  basicstyle=\ttfamily\small,
  breaklines=true,
  showstringspaces=false
}

\newtcolorbox{outputbox}{
  enhanced,
  breakable,
  colback=outputbg,
  colframe=outputborder,
  boxrule=0.5pt,
  arc=2pt,
  left=6pt,right=6pt,top=4pt,bottom=4pt,
  fonttitle=\bfseries,
  title=Sample Output
}

\titleformat{\section}{\large\bfseries}{\thesection}{1em}{}
\titleformat{\subsection}{\normalsize\bfseries}{\thesubsection}{1em}{}

\title{Sample Inputs and Outputs \\ Python Programming Question Bank \\ Units 3, 4, 5, 6}
\author{Abhishek Y. Sapkale}
\date{}

\begin{document}
\maketitle
\tableofcontents
\newpage

% ============================================================
\section{Unit 3: Strings}
% ============================================================

\subsection{Q13: Reverse a string without using built-in functions}
\begin{outputbox}
Enter a string: hello\\
Reversed string: olleh
\end{outputbox}

\subsection{Q14: Check if a string is a palindrome}
\begin{outputbox}
Enter a string: madam\\
madam is a palindrome
\end{outputbox}

\subsection{Q15: Count vowels and consonants in a string}
\begin{outputbox}
Enter a string: Abhishek\\
Vowels: 3\\
Consonants: 5
\end{outputbox}

\subsection{Q16: Find length of a string without len()}
\begin{outputbox}
Enter a string: python\\
Length of string: 6
\end{outputbox}

\subsection{Q17: Convert lowercase to uppercase and vice versa}
\begin{outputbox}
Enter a string: Hello World\\
Converted string: hELLO wORLD
\end{outputbox}

\subsection{Q18: Remove spaces from a string}
\begin{outputbox}
Enter a string: I am a student\\
String without spaces: Iamastudent
\end{outputbox}

\subsection{Q19: Count frequency of each character}
\begin{outputbox}
Enter a string: aabbc\\
'a' : 2\\
'b' : 2\\
'c' : 1
\end{outputbox}

\subsection{Q20: Count words in a sentence}
\begin{outputbox}
Enter a sentence: I love python programming\\
Number of words: 4
\end{outputbox}

\subsection{Q21: Find largest and smallest word in a string}
\begin{outputbox}
Enter a sentence: I am learning python programming\\
Largest word: programming\\
Smallest word: I
\end{outputbox}

% ============================================================
\section{Unit 4: Lists and Tuples}
% ============================================================

\subsection{Q13: Create a list and print all its elements using a loop}
\begin{outputbox}
10\\
20\\
30\\
40\\
50
\end{outputbox}

\subsection{Q14: Find the largest and smallest element in a list}
\begin{outputbox}
Largest element: 67\\
Smallest element: 2
\end{outputbox}

\subsection{Q15: Demonstrate list functions (append, insert, remove, pop)}
\begin{outputbox}
Original list: [1, 2, 3, 4, 5]\\
After append(6): [1, 2, 3, 4, 5, 6]\\
After insert(2, 100): [1, 2, 100, 3, 4, 5, 6]\\
After remove(100): [1, 2, 3, 4, 5, 6]\\
After pop(): [1, 2, 3, 4, 5]
\end{outputbox}

\subsection{Q16: Create a nested list and access specific elements}
\begin{outputbox}
Full nested list: [[1, 2, 3], [4, 5, 6], [7, 8, 9]]\\
Element at row 0, col 2: 3\\
Element at row 1, col 1: 5\\
Element at row 2, col 0: 7
\end{outputbox}

\subsection{Q17: List comprehension to generate squares of numbers from 1 to 10}
\begin{outputbox}
Squares from 1 to 10: [1, 4, 9, 16, 25, 36, 49, 64, 81, 100]
\end{outputbox}

\subsection{Q18: Check whether an element exists in a list using membership operators}
\begin{outputbox}
Enter a number to search: 30\\
30 is present in the list
\end{outputbox}

\subsection{Q19: Create a tuple and demonstrate tuple packing and unpacking}
\begin{outputbox}
Packed tuple: (1, 2, 3)\\
Unpacked values: 1 2 3
\end{outputbox}

\subsection{Q20: Compare two tuples and display whether they are equal or not}
\begin{outputbox}
Tuples are equal
\end{outputbox}

% ============================================================
\section{Unit 5: Sets and Dictionaries}
% ============================================================

\subsection{Q11: Create a set and perform add and remove operations}
\begin{outputbox}
Original set: \{1, 2, 3, 4, 5\}\\
After add(6): \{1, 2, 3, 4, 5, 6\}\\
After remove(3): \{1, 2, 4, 5, 6\}
\end{outputbox}

\subsection{Q12: Find the union and intersection of two sets}
\begin{outputbox}
Union: \{1, 2, 3, 4, 5, 6\}\\
Intersection: \{3, 4\}
\end{outputbox}

\subsection{Q13: Check whether an element is present in a set using membership operators}
\begin{outputbox}
Enter a number to search: 25\\
25 is not present in the set
\end{outputbox}

\subsection{Q14: Create a set using set comprehension (squares of numbers)}
\begin{outputbox}
Set of squares: \{1, 4, 36, 70, 9, 16, 49, 81, 100, 25, 64\}
\end{outputbox}
\textit{Note: Set elements may not display in sorted order since sets are unordered in Python.}

\subsection{Q15: Create a dictionary and display its elements}
\begin{outputbox}
name : Abhishek\\
roll : AD1156\\
course : AI \& DS
\end{outputbox}

\subsection{Q16: Access and update values in a dictionary}
\begin{outputbox}
Name: Abhishek\\
Updated dictionary: \{'name': 'Abhishek', 'roll': 'AD1156', 'course': 'AI \& Data Science'\}
\end{outputbox}

\subsection{Q17: Delete an element from a dictionary using pop() and del}
\begin{outputbox}
After pop(): \{'name': 'Abhishek', 'roll': 'AD1156'\}\\
After del: \{'name': 'Abhishek'\}
\end{outputbox}

\subsection{Q18: Create a dictionary using dictionary comprehension (number and its square)}
\begin{outputbox}
Dictionary of squares: \{1: 1, 2: 4, 3: 9, 4: 16, 5: 25, 6: 36, 7: 49, 8: 64, 9: 81, 10: 100\}
\end{outputbox}

% ============================================================
\section{Unit 6: Functions and Modules}
% ============================================================

\subsection{Q13: Create a function to find the factorial of a number}
\begin{outputbox}
Enter a number: 5\\
Factorial: 120
\end{outputbox}

\subsection{Q14: Check whether a number is prime or not}
\begin{outputbox}
Enter a number: 7\\
7 is a prime number
\end{outputbox}

\subsection{Q15: Function that accepts a list and returns the sum of all elements}
\begin{outputbox}
Sum of list: 100
\end{outputbox}

\subsection{Q16: Lambda function to find the square of a number}
\begin{outputbox}
Enter a number: 9\\
Square: 81
\end{outputbox}

\subsection{Q17: Demonstrate the use of default and keyword arguments}
\begin{outputbox}
Name: Abhishek\\
Course: AI \& Data Science\\
Name: Om\\
Course: Computer Science
\end{outputbox}

\subsection{Q18: Recursive function to calculate Fibonacci series up to n terms}
\begin{outputbox}
Enter number of terms: 8\\
0 1 1 2 3 5 8 13
\end{outputbox}

\subsection{Q19: Create a module with addition and subtraction functions, then use it}
\begin{outputbox}
Addition: 15\\
Subtraction: 5
\end{outputbox}

\subsection{Q20: Import math module to calculate square root, power, and factorial}
\begin{outputbox}
Enter a number: 4\\
Square root: 2.0\\
Power (num\textasciicircum2): 16.0\\
Factorial: 24
\end{outputbox}

\end{document}